% Ethics Considerations
% Required by ACM CCS 2026

We performed a stakeholder-based ethics analysis following the Menlo Report framework~\cite{dittrich2012menlo}, considering four principles:

\paragraph{Beneficence.}
Our research directly benefits platform operators and users by identifying a systematic capability-description gap that existing platform scanners entirely miss (0\% T6 recall).
The detection pipeline we propose could prevent future supply-chain attacks affecting hundreds of users.

\paragraph{Respect for Persons.}
All data was collected from publicly accessible APIs.
No user accounts were accessed, no private data was collected, and no skill was executed during analysis.
T6 malicious skill identifiers are pseudonymized throughout the paper to protect affected developers who may have been victims of account compromise.

\paragraph{Justice.}
We applied our analysis uniformly across all skills in both ecosystems without targeting individual developers.
Our framework classifies skills, not developers---the high inconsistency rate (90\%)
reflects a widespread ecosystem norm, not individual malfeasance.

\paragraph{Respect for Law and Public Interest.}
We notified both platform operators of our findings prior to submission (responsible disclosure).
Our crawler respected \texttt{robots.txt} and rate-limited requests to $\leq$0.5 req/s.
No laws or terms of service were violated during data collection.

\paragraph{Potential harms.}
Publishing detailed capability extraction patterns could theoretically help attackers craft harder-to-detect malicious skills.
We mitigate this by (a)~withholding the most specific evasion-relevant pattern details from the public artifact and (b)~sharing the full pattern bank only with platform operators and academic researchers under data use agreements.
